% ===================================================================
%  GAMBAR No. 14 — Dalan. --- Para Pedagang.
%
%  Traduction, faite en 2026, en javanais d'aujourd'hui, sur l'ido de
%  texto/io. Regles de la colonne : voir 10-gambar-01.tex. Une seule
%  numerotation.
%
%  QUATRE-VINGT-DIX-NEUF RENVOIS ET PRESQUE AUTANT DE METIERS : c'est
%  le tableau des noms d'agent, et le javanais y a une machine a lui
%  qui rend le francais et l'ido lourds a cote.
%
%   * « tukang » + ce qu'on fait : tukang cukur le coiffeur (23),
%     tukang sapu le balayeur (11), tukang jait le tailleur (62),
%     tukang payung le parapluitier (51), tukang gawe timah le
%     ferblantier (14), tukang cerobong le ramoneur (91), tukang
%     potret le photographe (20).
%   * « juru » + ce dont on repond : juru tulis, juru masak, juru
%     gambar, juru warta.
%   * « bakul » + ce qu'on vend : bakul kembang la bouquetiere (92),
%     bakul jeruk la marchande d'oranges (94), bakul buku le libraire
%     (70).
%
%  TROIS PREFIXES, ET LE CHOIX ENTRE EUX N'EST PAS LIBRE. « tukang »
%  dit la main, « juru » dit la charge, « bakul » dit la vente ; un
%  juru masak n'est pas un tukang masak, et un bakul buku n'est pas un
%  tukang buku. Le francais et l'ido derivent tous leurs metiers d'un
%  seul suffixe — -iste, -isto —, et perdent en route ce que le
%  javanais dit en un mot.
%
%  ET « bakul » EST LE MOT QUI TIENT LE LIVRET ENSEMBLE : il a deja
%  servi au tableau 3 pour la laitiere (24), et il servira au tableau
%  15 pour tout le marche. C'est le mot le plus employe de cette
%  colonne apres les mots-outils.
%
%  LE COMMERCE MODERNE, LUI, EMPRUNTE COMME AUX TROIS TABLEAUX
%  PRECEDENTS : konsul (9), optik (38), fotograf, mesin tik la machine
%  a ecrire, stenograf (35), litograf (72), redaksi (71), montir,
%  advokat (82), notaris (84). La regle du tableau 11 se verifie une
%  quatrieme fois : ce qui est arrive tout fait garde son nom.
%
%  MAIS L'ENTERREMENT (4) EST ENTIEREMENT JAVANAIS, ET C'EST LA
%  MEILLEURE RENCONTRE DU TABLEAU. layon le corps, kranda le
%  cercueil, layadan le convoi, sungkawa le deuil, kuburan. Java a
%  enterre ses morts bien avant que quiconque arrive, et le seul mot
%  du cortege que la colonne emprunte est « pandhita », pour le clerge
%  — parce que le clerge grave est catholique et que « pandhita »,
%  qui est sanskrit, sert en javanais a tous les hommes de religion
%  sans distinction. Un mot ancien peut couvrir une chose nouvelle
%  quand il est assez large.
% ===================================================================

%%K t14-apar-1 apar
\VUsaut{4.0mm}
\VUtitre{15.0pt}{60}{GAMBAR No. 14}
\VUsaut{3.0mm}

%%K t14-tit-1 sub
\VUpk{11.4pt}{40}{Dalan. --- Para Pedagang.}
\VUsaut{3.0mm}

%%K t14-01-1 p
1. --- Kancaku Ioannes, kang manggon ing desa, teka nglampahi telung
dina ing kulawargaku. \VUgras{Nganti} wektu iku dheweke durung tau duwe
kesempatan ndeleng kutha gedhe ; mula kita nglampahi kabeh dinane karo
mlaku-mlaku, lan aku nerangake marang dheweke apa kang durung
diweruhi.

%%K t14-02-1 p
2. --- Kang kapisan kita lunga nyowani \VUgras{observatorium}
\textsuperscript{(1)}, kang narik banget kawigatene. Nalika bali liwat
\VUgras{dalan} \textsuperscript{(3)} panggonan \VUgras{papan adus Turki}
\textsuperscript{(2)}, kita kepethuk \VUgras{layadan}
\textsuperscript{(4)}. Ing ngarepe \VUgras{iring-iringan layadan} mlaku
\VUgras{para pandhita,} ndhisiki \VUgras{kreta layon,} panggonan
\VUgras{kranda} diselehake. Akeh kanca ngeterake kulawarga kang lagi
\VUgras{sungkawa.}

%%K t14-02-2 p
Sawise iring-iringan liwat, kita nyabrang dalan ing ngarepe
\VUgras{warung bir} \textsuperscript{(5)} gedhe panggonan akeh \VUgras{tamu}
ngombe \VUgras{bir} ing teras, kaeyuban \VUgras{kanopi kaca}
\textsuperscript{(6)}.

%%K t14-03-1 p
3. --- Ioannes gumun banget marang \VUgras{lambang} kang dipasang ing
antarane tokone \VUgras{bakul ciu} \textsuperscript{(7)} lan tokone
\VUgras{bakul musik} \textsuperscript{(10)}. Dheweke takon marang aku apa
tegese ; aku mangsuli yen iku nuduhake \VUgras{kantor konsul}
\textsuperscript{(8)} Inggris, lan aku nuduhake marang dheweke
\VUgras{konsul} \textsuperscript{(9)} kang ana ing balkon.

%%K t14-tit-2 sub
\VUpk{10.2pt}{40}{Ndelok-ndelok toko.}
\VUsaut{3.0mm}

%%K t14-04-1 p
4. --- Mesthi wae kita mandheg ing ngarepe pajangan \VUgras{alat musik,}
\VUgras{harmonium,} \VUgras{orgel} lan \VUgras{piano} ; kita nyawang
tutup buku \VUgras{partitur} lan \VUgras{lagu} kanggo piano kang ana
gambare, lan gumun marang \VUgras{kecapi} kayu kang disepuh emas, kang
dianggo dadi papan jeneng.

%%K t14-05-1 p
5. --- Mendhung lebu kang digawe \VUgras{tukang sapu} \textsuperscript{(11)}
lan \VUgras{kreta sapu} \textsuperscript{(12)} meksa kita enggal liwat ing
ngarepe \VUgras{perabot} endah duweke \VUgras{bakul perabot}
\textsuperscript{(13)}, lan ing ngarepe pajangan \VUgras{tungku} lan
\VUgras{lampu} duweke \VUgras{tukang gawe timah} \textsuperscript{(14)}.

%%K t14-06-1 p
6. --- Kejaba iku, ing panggonan kono, kita kepeksa ninggal trotoar,
amarga trotoare kebak buntelan kang lagi diudhunake saka \VUgras{kreta
angkutan} \textsuperscript{(16)} kanggo dilebokake ing \VUgras{gudhang}
\textsuperscript{(15)} kang lagi wae disewa.

%%K t14-06-2 p
Iku ngalang-alangi kita weruh kreta endah duweke \VUgras{tukang gawe
kreta} \textsuperscript{(17)}, nanging ora ngalang-alangi kita mandheg
nyawang \VUgras{tukang pak} \textsuperscript{(19)} kang lagi gawe peti
kanggo tanggane, yaiku \VUgras{bakul pit lan montor}
\textsuperscript{(18)}. Banjur, sarehne wis rada kasep, kita bali menyang
omah.

%%K t14-tit-3 sub
\VUpk{10.2pt}{40}{Mlaku-mlaku ing kutha.}
\VUsaut{3.0mm}

%%K t14-07-1 p
7. --- Ing dina esuke, ibuku ngusulake marang Ioannes supaya dipotret,
banjur ibuku ngutus aku ngeterake dheweke menyang \VUgras{tukang
potret} \textsuperscript{(20)}. Sawise sarapan awan, kita mlebu ing
\VUgras{papan gawe} sang seniman, panggonan kita kepeksa ngenteni rada
suwe. Kanggo nglong wektu, kita nyawang \VUgras{potret}
\textsuperscript{(22)} kang gumantung ing tembok.

%%K t14-07-2 p
Nalika giliran kita, pagaweane ora suwe, lan sanalika kancaku rampung
mapan ing ngarepe \VUgras{kamera} \textsuperscript{(21)}, kita mudhun.

%%K t14-08-1 p
8. --- Ing lantai kapisan, kita weruh liwat lawange \VUgras{tukang
cukur} \textsuperscript{(23)}, kang menga, kancane lagi ngethok rambute
tamu.

%%K t14-08-2 p
Tekan ing dalan, kita liwat ing ngarepe \VUgras{toko rokok}
\textsuperscript{(24)}. Ioannes seneng banget marang \VUgras{lentera}
\textsuperscript{(25)} abang lan \VUgras{pipa gedhe} kang dianggo
\VUgras{papan jeneng} \textsuperscript{(26)}, nganti dheweke nubruk wong
tuwa loro kang lagi ngobrol karo \VUgras{nyerot bako}
\textsuperscript{(27)}.

%%K t14-tit-4 sub
\VUpk{10.2pt}{40}{Toko roti manis.}
\VUsaut{3.0mm}

%%K t14-09-1 p
9. --- \VUgras{Dhuwit kertas} lan \VUgras{dhuwit receh} saka sakabehe
negara kang dipajang ing \VUgras{etalase} \VUgras{tukang ijol dhuwit}
\textsuperscript{(29)} narik kawigaten kita sedhela, nanging sarehne wis
jame mangan cilik, kita mlebu menyang \VUgras{tukang gawe roti manis}
\textsuperscript{(31)} tanpa mandheg suwe-suwe.

%%K t14-10-1 p
10. --- Weruh kabeh \VUgras{panganan legi} kang ana ing toko iku —
\VUgras{roti tar, roti madu, biskuit} lan warna-warna \VUgras{permen} —
kita angel milih. Pungkasane Ioannes milih \VUgras{roti krim,} lan aku
mung njupuk \VUgras{tart ceri} cilik, amarga panganan gula \VUgras{gawe
lara untuku,} lan aku babar pisan ora kepengin kerep-kerep nyowani
\VUgras{dhokter untu} \textsuperscript{(30)}.

%%K t14-tit-5 sub
\VUpk{10.2pt}{40}{Nulis nganggo mesin.}
\VUsaut{3.0mm}

%%K t14-11-1 p
11. --- Kanggo nuduhake marang kancaku \VUgras{mesin tik} kang wis tau
dirungu nanging durung diweruhi, kita mlebu \VUgras{kantor}
\textsuperscript{(32)} \VUgras{sedulur lanangku} \textsuperscript{(34)}.
Kita nemu dheweke lagi ndikte layange marang \VUgras{sekretaris}e
\textsuperscript{(35)}, kang dadi \VUgras{juru stenografi.} Lagi wae
layange rampung, \VUgras{juru ketik} \textsuperscript{(33)} nyalin nganggo
mesine. Sarehne ora kepengin nyelani pagaweane sanakku, kita mung
sedhela ana ing kono.

%%K t14-12-1 p
12. --- Ing sangisore kantore sedulurku, manggon \VUgras{tukang jam lan
mas-masan} \textsuperscript{(37)} kang gedhe, kang uga tukang gawe emas.
Tokone kang endah banget isi akeh \VUgras{jam, jam bandhul, jam
kanthong, ranté cilik} lan \VUgras{mas-masan} kang larang, tuladhane
\VUgras{kalung mutiara,} \VUgras{gelang} mas, \VUgras{anting} lan
\VUgras{ali-ali} kang direngga \VUgras{sesotya} lan \VUgras{inten.}
Salah sijine etalase mligi disedhiyakake kanggo \VUgras{barang emas}
lan isi klumpukan gedhe \VUgras{piranti mangan} salaka.

%%K t14-13-1 p
13. --- Nalika menggok ing pojoke dalan, aku weruh yen \VUgras{papan
jeneng} \textsuperscript{(36)} kang dipasang cedhak \VUgras{nomer} omah
wis diganti. Aku arep ngucapake pengamatanku kanthi seru, nalika swara
bungahe \VUgras{musik} prajurit krungu. Kita mlayu mapag resimen, tanpa
mikir yen \VUgras{resimen} iku bakal liwat ing ngarepe tokone
\VUgras{tukang topi} \textsuperscript{(39)} lan \VUgras{tukang sarung
tangan} \textsuperscript{(40)}, lan yen kita bakal padha becike weruh
arak-arakane yen tetep ngadeg ing ngarepe etalase \VUgras{tukang
optik} \textsuperscript{(38)}, kang kebak \VUgras{kacamata cepit,
kacamata gagang} lan \VUgras{teropong cilik.}

%%K t14-tit-6 sub
\VUpk{10.2pt}{40}{Arak-arakan mlaku.}
\VUsaut{3.0mm}

%%K t14-14-1 p
14. --- Ing ngarepe \VUgras{resimen} \textsuperscript{(41)} mlaku
\VUgras{kepala tambur} \textsuperscript{(42)} kang dietutake \VUgras{para
tukang tambur} \textsuperscript{(43)}, \VUgras{para tukang trompet}
\textsuperscript{(44)} lan kabeh \VUgras{pemain musik.}
\VUgras{Jendral} \textsuperscript{(45)} kang ana ing alun-alun bareng karo
ajudane, mandheg cedhak \VUgras{pancuran} \textsuperscript{(49)}
\VUgras{sendhang} \textsuperscript{(48)} kanggo nyekseni
\VUgras{arak-arakan.}

%%K t14-14-2 p
\VUgras{Para opsir staf} \textsuperscript{(46)}, \VUgras{kolonel} lan
\VUgras{letnan kolonel} aweh salam nganggo pedhang. \VUgras{Para
mayor} ana ing ngarepe \VUgras{batalion}e \textsuperscript{(47)}, lan
saben \VUgras{kapten} mrentah \VUgras{kompi}ne, dene \VUgras{para
letnan, letnan enom} lan \VUgras{bintara} manggon ing panggonane kang
lumrah ing sacedhake anak buahe.

%%K t14-15-1 p
15. --- Akeh wong liwat padha nglumpuk ing \VUgras{prapatan}
\textsuperscript{(50)} ; para pedagang, \VUgras{tukang payung}
\textsuperscript{(51)}, \VUgras{tukang gawe barang timah}
\textsuperscript{(53)}, \VUgras{tukang gegaman} \textsuperscript{(54)}
padha metu kanggo weruh \VUgras{para prajurit.} Kabeh kreta —
\VUgras{fiakre} \textsuperscript{(52)}, \VUgras{kreta tiang}
\textsuperscript{(57)} lan uga \VUgras{tong panyiram}
\textsuperscript{(56)} — padha minggir sadawane trotoar.

%%K t14-tit-7 sub
\VUpk{10.2pt}{40}{Adegan ing dalan.}
\VUsaut{3.0mm}

%%K t14-15-2 p
\VUgras{Salesman kang mubeng} \textsuperscript{(55)}, kang nyangking
\VUgras{kothak conto}ne, enggal nyabrang dalan, lan wong-wong kang
lungguh ing \VUgras{lantai ndhuwur} \textsuperscript{(59)} \VUgras{omnibus}
\textsuperscript{(58)} padha noleh kanggo ngrasakake tontonan iku.
\VUgras{Penyanyi keliling} \textsuperscript{(60)} kang mesakake, karo
\VUgras{orgel puteran}e \textsuperscript{(61)}, kepeksa nyelani
\VUgras{lagu}ne, amarga swarane ketutupan bunyine tambur.

%%K t14-16-1 p
16. --- Nalika resimen tekan ing ngarepe \VUgras{toko kain}
\textsuperscript{(64)}, \VUgras{para tukang jait wadon}
\textsuperscript{(63)} lan \VUgras{para tukang jait} \textsuperscript{(62)}
ninggal \VUgras{mesin jait}e lan pagaweane kanggo nyawang liwat
jendhela. \VUgras{Sing duwe toko} \textsuperscript{(65)}, kang lagi wae
masang \VUgras{sandhangan} \textsuperscript{(66)} anyar ing
\VUgras{pajangan}e, kepeksa nglebokake \VUgras{kain}
\textsuperscript{(67)} lan \VUgras{lawon} kang dipasang ing trotoar,
cedhak \VUgras{bolongan got} \textsuperscript{(68)}, wedi yen wong akeh
bakal nabrak ; malah \VUgras{bakul mider} \textsuperscript{(69)} tuwa,
kang lagi ditawar kaos kaki dening penjaga lawang, kepeksa mlebu
lorong kanggo ngrampungake dodolane.

%%K t14-16-2 p
Kita ngeterake resimen tekan tangsi, \VUgras{lan,} kanthi marem banget
marang dinane, kita bali ing jame bujana bengi.

%%K t14-tit-8 sub
\VUpk{10.2pt}{40}{Werna-werna dagang lan pagawean.}
\VUsaut{3.0mm}

%%K t14-17-1 p
17. --- Ing dina sabanjure, bapakku ngusulake kita nyowani percetakane
koran kono. Kita nampa kanthi bungah.

%%K t14-17-2 p
\VUgras{Tukang cetak} iku uga \VUgras{bakul buku} \textsuperscript{(70)},
\VUgras{penerbit,} \VUgras{bakul kertas} lan \VUgras{tukang jilid.}
Dheweke masang ing lantai kapisan \VUgras{kantor redaksi}
\textsuperscript{(71)} koran, lan uga \VUgras{papan ukir,} panggonan
\VUgras{tukang litografi} \textsuperscript{(72)} lan \VUgras{tukang ukir
tembaga} \textsuperscript{(73)} nyambut gawe, sarta pungkasane
\VUgras{papan nyusun huruf,} panggonan \VUgras{buruh cetak}
\textsuperscript{(74)} ana. \VUgras{Percetakan} \textsuperscript{(75)} kang
sanyatane, \VUgras{mesin cetak} lan warna-warna mesin ana ing lantai
ngisor.

%%K t14-tit-9 sub
\VUpk{10.2pt}{40}{Ioannes akeh takone.}
\VUsaut{3.0mm}

%%K t14-18-1 p
18. --- Nalika metu saka percetakan, Ioannes gumun banget lan mandheg
ing ngarepe kothak kaca cilik kang dipasang ing cagak, adhep-adhepan
karo \VUgras{toko benang lan dom} \textsuperscript{(77)}, banjur takon
kanggo apa kothak iku. Bapakku nerangake marang dheweke yen iku
\VUgras{alarm kobongan} \textsuperscript{(76)}. Kejaba iku kabeh, ing
kutha, gawe gumune kancaku, wiwit \VUgras{sendhang watu}
\textsuperscript{(78)} kang prasaja lan \VUgras{sepedha motor rodha telu
kanggo momotan} \textsuperscript{(80)} kang entheng, kang dituntun
\VUgras{bocah kurir} \textsuperscript{(79)} kang nganggo seragam, nganti
\VUgras{kreta pos} \textsuperscript{(81)}.

%%K t14-19-1 p
19. --- Nalika bapakku sedhela ninggal kita kanggo ngomong karo
\VUgras{juru narik pajeg} \textsuperscript{(83)}, kang mandheg kanggo
ngobrol karo \VUgras{advokat} \textsuperscript{(82)} lan \VUgras{notaris}
\textsuperscript{(84)}, kita seneng-seneng ngetutake nganggo mripat
mlaku-mlakune wong ing dalan. Wong kang maneka warna liwat ing ngarepe
kita : \VUgras{kancane tukang roti manis} \textsuperscript{(85)}, kang
nggawa \VUgras{pastel} kang endah ing kranjang, teka ing mburine
\VUgras{bakul sandhangan} \textsuperscript{(86)} ; \VUgras{tukang urub
lentera} \textsuperscript{(87)} lagi ngobrol karo \VUgras{tukang gas}
\textsuperscript{(88)} ; \VUgras{suster} \textsuperscript{(89)} kang nyekel
tangane bocah lola wadon lagi bali menyang biarane.

%%K t14-tit-10 sub
\VUpk{10.2pt}{40}{Nalika kita bali menyang omah.}
\VUsaut{3.0mm}

%%K t14-20-1 p
20. --- Pas nalika bapakku nyusul kita, Ioannes ngguyu kepingkel-pingkel
weruh rai reged lan sandhangan aneh loro \VUgras{tukang cerobong}
\textsuperscript{(91)} kang kebak langes.

%%K t14-21-1 p
21. --- Aku kepengin tuku tanduran saka \VUgras{bakul kembang}
\textsuperscript{(92)} kanggo ibuku, nanging bapakku ngelingake yen
tandurane bakal abot banget digawa lan luwih becik tuku \VUgras{jeruk.}
Kita nyedhaki \VUgras{bakul} \textsuperscript{(94)} iku, lan bareng
\VUgras{pengasuh} \textsuperscript{(93)}, kang lagi milihake kanggo bocah
asuhane, rampung tukune, kita dilayani.

%%K t14-22-1 p
22. --- Nalika bali menyang omah, kita kepethuk sawetara wong kang
tepung : kancane kulawargaku kang wis tuwa, kang nyekeli lengene
\VUgras{nyonyah kancane} \textsuperscript{(95)}, \VUgras{insinyur}
\textsuperscript{(96)} kreteg lan dalan, lan salah sijine sedulurku wadon
karo \VUgras{pengasuh bocah}e \textsuperscript{(98)}.

%%K t14-22-2 p
Banjur kita kepethuk \VUgras{tukang topi wadon} \textsuperscript{(97)}
ibuku lan \VUgras{biarawan} \textsuperscript{(99)} loro kang manca ing
kutha kono, kang nyedhaki kita kanggo takon marang bapakku bab dalan
menyang setasiun.

\VUsaut{6.0mm}
\VUfilet{22mm}
